\section{方法设计}
\subsection{总体框架}
本文提出的 \method{} 建立在“冻结编码器 + 特征级可信输出层”的基本思想之上。给定输入文本 $x_i$，首先利用长上下文编码器提取 token 级隐藏状态，并通过池化操作获得定长表示 $\mathbf{z}_i$；随后在输出层中同时预测标签均值 $\boldsymbol{\mu}_i$ 与输入相关方差 $\boldsymbol{\sigma}_i^2$；最后根据方差构造样本权重，对均值分支的最后一层进行加权最小二乘精修。经过这一过程，模型输出不再只是单一得分，而是由“预测值 + 噪声估计 + 结构化修正”共同组成。

与传统做法相比，该框架的核心特点在于将表示学习与最终映射求解适度解耦。编码器负责提供稳定的语义表示，输出层则专注于对预测不确定性和样本差异进行建模。这样的分工使本文能够在不大幅增加训练成本的前提下，把研究重点集中到输出层的可信性设计上。

\subsection{特征表示与池化}
设编码器最后一层隐藏状态为 $\mathbf{H}_i \in \mathbb{R}^{L_i \times d}$，注意力掩码为 $\mathbf{m}_i \in \{0,1\}^{L_i}$，则本文采用掩码均值池化构造样本表示：
\[
\mathbf{z}_i = \frac{\sum_{t=1}^{L_i} m_{it} \mathbf{h}_{it}}{\sum_{t=1}^{L_i} m_{it}}.
\]
均值池化的优点在于形式简单、参数开销低，并能在长文本中整合分散出现的证据。对于临床场景而言，许多标签并不依赖单个显著触发词，而是由多个局部线索共同支撑，因此基于整段文本的平均表示是一种稳健的起点。

当然，均值池化并不意味着所有位置的重要性完全一致。它更像是一种保证训练稳定性的基础聚合方式。后续若需要进一步强化关键句建模，可在同一框架下替换为注意力池化或其他更细粒度的汇聚策略，但这并不影响本文关于可信输出层的核心讨论。

\subsection{基线输出层}
为了更清晰地说明本文方法的增益来源，首先考虑三类常见输出层。线性头直接学习
\[
\hat{\mathbf{y}}_i = \mathbf{W}\mathbf{z}_i + \mathbf{b},
\]
对应最基本的特征到标签映射。MLP 头在此基础上加入隐藏层和非线性激活，以增强特征变换能力。Ridge 头则进一步在目标函数中加入 $L_2$ 正则项：
\[
\hat{\boldsymbol{\beta}} = \arg\min_{\boldsymbol{\beta}} \norm{\mathbf{X}\boldsymbol{\beta} - \mathbf{Y}}_2^2 + \lambda \norm{\boldsymbol{\beta}}_2^2.
\]
上述三类基线分别对应“简单线性映射”“浅层非线性映射”和“带显式正则约束的线性映射”。

\method{} 与这些基线的本质区别，不在于网络更深，而在于输出层显式表示样本相关噪声，并利用该噪声进一步修正最终参数。因此，本文将这些方法作为对照，以检验输出层可信建模是否具有独立价值。

\subsection{双分支异方差建模}
在本文框架中，输出层由均值分支与方差分支组成。给定特征 $\mathbf{z}_i$，均值分支产生
\[
\boldsymbol{\mu}_i = g_{\mu}(\mathbf{z}_i),
\]
方差分支产生
\[
\log \boldsymbol{\sigma}_i^2 = g_{\sigma}(\mathbf{z}_i).
\]
随后通过指数映射和稳定常数 $\epsilon$ 得到正的方差估计：
\[
\boldsymbol{\sigma}_i^2 = \exp(\log \boldsymbol{\sigma}_i^2) + \epsilon.
\]
进一步地，将其转化为样本权重
\[
\mathbf{w}_i = \frac{1}{\boldsymbol{\sigma}_i^2}.
\]
其中，较小的预测方差意味着该样本在相应输出维度上更可信，从而拥有更大的权重；较大的预测方差则意味着证据不足、噪声较强或标签边界不清，对参数估计的影响应适当减弱。

这一设计的直观动机来自临床文本的异质性。不同病历的证据密度、语言规范性和标签清晰度并不相同，若所有样本在训练中被同等对待，模型就难以显式区分容易样本与高风险样本。通过异方差建模，输出层能够学习到与样本难度相关的噪声结构，从而为后续精修提供依据。

\subsection{异方差训练目标}
若从概率建模角度出发，异方差高斯负对数似然可写为
\[
\mathcal{L}_{\text{NLL}} = \frac{1}{2}\left(\log \sigma^2 + \frac{(y-\mu)^2}{\sigma^2}\right).
\]
这一目标能够同时约束均值预测与方差估计，但在实践中容易出现方差分支过度放大、训练不稳定等问题。为缓解这一现象，本文采用 Beta-NLL 形式的稳定化目标：
\[
\mathcal{L}_{\beta\text{-NLL}} = \frac{\mathcal{L}_{\text{NLL}}}{(\sigma^2_{\text{detach}})^{\beta}},
\]
其中 $\beta$ 用于控制重加权强度，detach 操作避免方差项直接主导均值分支梯度。

该目标的作用不是单纯追求概率意义上的完备性，而是在保留异方差建模能力的同时，提高训练过程的稳定性。对于本文关注的输出层研究而言，稳定训练比过于理想化的概率解释更重要，因为只有在训练可控的前提下，方差估计才能进一步服务于结构化精修。

\subsection{阶段 A：联合学习均值与方差}
在第一阶段，本文对均值分支和方差分支进行联合训练。此时模型通过梯度下降学习从冻结特征到预测均值、预测方差的映射关系。经过这一阶段后，模型已经具备基础的不确定性表达能力，并能够给出样本级噪声估计。

这一阶段的作用可以理解为可信输出层的预估计过程。均值分支负责形成初始预测，方差分支则学习哪些样本在当前任务中更稳定、哪些样本更容易产生误差。尽管阶段 A 本身已经能够带来性能增益，但它仍然完全依赖神经网络参数化表达，尚未充分利用显式统计求解器的优势。

\subsection{阶段 B：基于样本权重的 WLS 精修}
在第二阶段，本文固定方差分支，仅利用其产生的样本权重对均值分支最后一层进行加权最小二乘精修。设最后一层前的隐藏表示为 $\mathbf{h}_i$，则对第 $j$ 个输出维度求解
\[
\hat{\boldsymbol{\beta}}_j = \arg\min_{\boldsymbol{\beta}_j} \sum_{i=1}^{N} w_{ij}\left(y_{ij} - \mathbf{h}_i^{\top}\boldsymbol{\beta}_j\right)^2 + \lambda \norm{\boldsymbol{\beta}_j}_2^2.
\]
其正规方程可写为
\[
(\mathbf{H}^{\top}\mathbf{W}_j\mathbf{H} + \lambda \mathbf{I})\boldsymbol{\beta}_j = \mathbf{H}^{\top}\mathbf{W}_j\mathbf{y}_j.
\]
求解完成后，以解析方式更新最后一层参数，从而得到新的输出映射。

这一过程可以看作对神经网络最终判别边界的全局重估。前端网络负责学习复杂特征和噪声模式，WLS 则在此基础上利用样本权重重新平衡不同样本对最后一层参数的贡献。与纯梯度下降相比，解析求解能够更直接地把“哪些样本更可信”这一信息体现在最终参数中，因此构成了本文方法区别于普通 MLP 头的关键一步。

\subsection{数值稳定性与方法边界}
由于 WLS 精修涉及线性系统求解，其可靠性取决于特征矩阵的数值性质和正则化强度。为保证方法可用，求解阶段需要避免直接矩阵求逆，并对病态情况进行必要监控；同时，正则系数 $\lambda$ 也需在防止过拟合与避免欠拟合之间取得平衡。换言之，结构化输出层并不是“加一个解析解就一定更好”，它依赖合理的数值稳定性设计。

此外，本文当前主要建模的是输入相关的偶然不确定性，而非完整的认知不确定性。预测方差可以为样本难度和输出稳健性提供重要线索，但它并不自动等同于严格校准后的置信度。因此，本文方法更适合作为临床可信预测的第一步：先在输出层建立噪声感知能力，再进一步扩展到校准、拒识和贝叶斯化分析。
